\subsection{General rod reduction and force accounting}
\label{sm:general}
For a spatial rod, define reference strains and body angular velocity by
\begin{equation}
 \Gamma=D^{\mathsf T}r_s,\qquad
 D_s=D\widehat\kappa,\qquad
 D_t=D\widehat\omega_b.
 \label{sm:rod-kinematics}
\end{equation}
The hat map satisfies $\widehat a b=a\times b$.
The spatial balance laws are
\begin{align}
 \rho A_c r_{tt}&=n_s+f_{\rm mag}+f_{\rm h}+f_{\rm c}+\rho A_c g,\\
 \partial_t(D\mathcal J_b\omega_b)&=m_s+r_s\times n+\tau_{\rm ext}.
 \label{sm:rod-balances}
\end{align}
Here $\mathcal J_b$ is the rotational inertia per unit reference
length. For constant body inertia, the second left-hand side equals
$D(\mathcal J_b\dot\omega_b+\omega_b\times\mathcal J_b\omega_b)$.
A stored-energy density determines the unconstrained force and moment
components through its strain derivatives \citep{simo,antman}.
In an unshearable specialization, the eliminated shear strains have
constraint reactions; setting those strains to zero in an unconstrained
shear law does not set the transverse internal forces to zero.

A body-fixed moment density $m_b(s)$ in a static applied field has
potential
\begin{equation}
 V_{\rm mag}=-\int_0^L (Dm_b)\cdot B_0(r)\dd s.
 \label{sm:distributed-magnetism}
\end{equation}
The point-dipole benchmark replaces this distribution by a distal
localized core. The permanent-dipole formulas are the leading terms of
a localized magnetic stress resultant \citep{abbott}. In a surrounding
vacuum region, the Maxwell stress is
\begin{equation}
 T_M=\mu_0^{-1}\left(BB^{\mathsf T}-\tfrac12|B|^2I\right).
 \label{sm:maxwell-stress}
\end{equation}
Its surface force and moment use the total field. Equivalently, for a
complete localized steady current $j(\xi)$, including bound currents,
and externally generated $B_{\rm ext}$,
\begin{align}
 m_d&=\tfrac12\int \xi\times j(\xi)\dd V,\nonumber\\
 F&=\nabla_r[m_d\cdot B_{\rm ext}(r)]+R_F,\nonumber\\
 \tau&=m_d\times B_{\rm ext}(r)+R_\tau.
 \label{sm:dipole-expansion}
\end{align}
Current closure gives $\int j\dd V=0$ and
$\int\xi_i j_j\dd V=\epsilon_{ijk}(m_d)_k$.
Taylor expansion then yields, with
$J_2=\int|\xi|^2|j(\xi)|\dd V$,
\begin{equation}
 |R_F|\le\tfrac12J_2\sup\|D^2B_{\rm ext}\|,\qquad
 |R_\tau|\le J_2\sup\|DB_{\rm ext}\|.
 \label{sm:dipole-remainders}
\end{equation}
The suprema cover the core and the Taylor segments. The self-field of
the complete isolated element has no net self-force or self-torque.
Thus one uses either the stress resultant or its dipole approximation,
without counting both as separate loads. The approximation
$\mu\simeq B_ra_m^3/\mu_0$ also assumes an effectively uniform,
body-fixed magnetization; a grade specification is not a measured
moment of an assembled tip.

Fluid and contact terms require their own constitutive choices. For
example, a local added-mass and drag closure may be written as
\begin{align}
 f_{\rm h}&=m_f a_f+m_a(a_f-r_{tt})
                  -c_1v_r-c_2|v_r|v_r,\nonumber\\
 v_r&=r_t-v_f,\qquad c_1,c_2\ge0.
 \label{sm:fluid-example}
\end{align}
This is a reduced closure to be calibrated, not a resolved fluid
interaction. In still fluid the drag power is nonpositive. In moving
fluid it is dissipative relative to $v_f$ and may supply energy to the
rod. A consistently defined added mass may be incorporated into the
kinetic metric once; it must then be removed from the external inertial
load. For a stationary wall, penetration $\delta>0$ and
$n_c=\nabla\delta$ give one nonadhesive penalty example:
\begin{align}
 V_c&=\frac{k_c}{p_c+1}[\delta]_+^{p_c+1},\nonumber\\
 f_c&=-k_c[\delta]_+^{p_c}n_c
       -c_c{\bf1}_{\{\delta>0\}}[\dot\delta]_+n_c.
 \label{sm:contact-example}
\end{align}
The second term dissipates energy during compression. Such switching
contact laws are excluded from the analytic neighbourhood used for the
persistence results. Reduced gravity alone does not remove fluid
forcing or select a blood-flow law. Neither
\eqref{sm:fluid-example} nor \eqref{sm:contact-example} is activated in
the revised numerical benchmark.

For a configuration map satisfying the constraints, define
\begin{equation}
 r_i=\partial_{q_i}r,\qquad
 \psi_i=\operatorname{vee}(D^{\mathsf T}\partial_{q_i}D).
 \label{sm:configuration-derivatives}
\end{equation}
Then $\omega_b=\sum_i\psi_i\dot q_i$, and
\begin{align}
 M_{ij}(q)&=\int_0^L
     [\rho A_c\,r_i\cdot r_j+\psi_i^{\mathsf T}\mathcal J_b\psi_j]\dd s,
       \nonumber\\
 Q_i&=\int_0^L[f\cdot r_i+\tau\cdot D\psi_i]\dd s+Q_i^\partial.
 \label{sm:galerkin-metric}
\end{align}
Concentrated tip masses and their moments of inertia add their
corresponding kinetic terms. A centreline interpolation by itself does
not specify rotational inertia.

Variation of $T-V$ gives
\begin{align}
 (c_{\rm geo})_i&=\sum_{j,k}\Gamma_{ijk}\dot q_j\dot q_k,\nonumber\\
 \Gamma_{ijk}&=\tfrac12
   (\partial_{q_j}M_{ik}+\partial_{q_k}M_{ij}
                         -\partial_{q_i}M_{jk}).
 \label{sm:christoffel}
\end{align}
Writing $p=M(q)\dot q$ gives
\begin{align}
 \dot p_i&=-\tfrac12p^{\mathsf T}(\partial_{q_i}M^{-1})p
       -\partial_{q_i}V\nonumber\\
 &\quad-\partial_{v_i}\mathcal D_N+(Q_{\rm drv}+B_u u)_i.
 \label{sm:variable-mass-momentum}
\end{align}
If the dissipation is stored as a momentum readout
$G_1(q,p)=\mathcal D_N(q,M^{-1}p)$, its physical force is
\begin{equation}
 -\partial_v\mathcal D_N=-M(q)\nabla_pG_1,
 \label{sm:rayleigh-momentum}
\end{equation}
not generally $-\nabla_pG_1$. The Christoffel term already follows from
$H_q$ and is not placed in $G_1$.
The energy identity is
\begin{equation}
 \dot H=-v^{\mathsf T}\partial_v\mathcal D_N
              +v^{\mathsf T}(Q_{\rm drv}+B_u u).
 \label{sm:general-power}
\end{equation}
Explicit time dependence adds $\partial_tH$; a time-dependent mass
also generates the corresponding inertial terms before the Legendre
transform. A quadratic positive semidefinite Rayleigh potential, or a
differentiable convex velocity potential minimized at zero, ensures
$v^{\mathsf T}\partial_v\mathcal D_N\ge0$. Nonnegativity of the potential
alone would not ensure that inequality.

\subsection{Weak form, endpoint inertia and consistent Hermite matrices}
\label{sm:weak-form}
For an admissible test displacement $\eta$ with
$\eta(0)=\eta_s(0)=0$, the benchmark weak form is
\begin{align}
 &\int_0^L\!\rho A_c w_{tt}\eta\dd s
       +m_tw_{tt}(L)\eta(L)+J_t\ddot\vartheta\eta_s(L)\nonumber\\
 &+\alpha\left[
       \int_0^L\!\rho A_c w_t\eta\dd s
       +m_tw_t(L)\eta(L)+J_t\dot\vartheta\eta_s(L)\right]\nonumber\\
 &+\int_0^L\!EIw_{ss}\eta_{ss}\dd s
       =F\eta(L)+[\tau_B-\mu B_0\sin\vartheta]\eta_s(L).
 \label{sm:beam-weak}
\end{align}
Two spatial integrations by parts give the interior and endpoint
equations in \eqref{mech:beam}--\eqref{mech:endpoint-balance}.
The $J_t$ term is a tip rigid-body inertia, distinct from the
distributed rotary inertia omitted by the Euler--Bernoulli model.

For an element of length $\ell_e$ and local coordinate
$\xi\in[0,1]$, the displacement row is
\begin{equation}
 N_e(\xi)^{\mathsf T}=
 \begin{pmatrix}
 1-3\xi^2+2\xi^3\\
 \ell_e(\xi-2\xi^2+\xi^3)\\
 3\xi^2-2\xi^3\\
 \ell_e(-\xi^2+\xi^3)
 \end{pmatrix}.
 \label{sm:hermite-shape}
\end{equation}
The local order is $(w_i,\vartheta_i,w_{i+1},\vartheta_{i+1})$.
Integrating the bending and kinetic energies yields
\begin{align}
 K_e&=\frac{EI}{\ell_e^3}
 \begin{pmatrix}
 12&6\ell_e&-12&6\ell_e\\
 6\ell_e&4\ell_e^2&-6\ell_e&2\ell_e^2\\
 -12&-6\ell_e&12&-6\ell_e\\
 6\ell_e&2\ell_e^2&-6\ell_e&4\ell_e^2
 \end{pmatrix},\nonumber\\
 M_e&=\frac{\rho A_c\ell_e}{420}
 \begin{pmatrix}
 156&22\ell_e&54&-13\ell_e\\
 22\ell_e&4\ell_e^2&13\ell_e&-3\ell_e^2\\
 54&13\ell_e&156&-22\ell_e\\
 -13\ell_e&-3\ell_e^2&-22\ell_e&4\ell_e^2
 \end{pmatrix}.
 \label{sm:element-matrices}
\end{align}
After assembly and elimination of the two clamped degrees of freedom,
add $m_t$ to the terminal displacement mass entry and $J_t$ to the
terminal rotation entry. Add $\mu B_0$ to the terminal rotation
stiffness entry only when forming the biased linear eigenproblem.
The nonlinear physical energy still uses $K_{\rm beam}$ and the cosine
potential.

For an axial centre-of-mass offset $e$ from the endpoint, the rigid-body
kinetic contribution would instead contain the block
\begin{equation}
 \begin{pmatrix}m_t&m_te\\m_te&J_{\rm cm}+m_te^2\end{pmatrix}
 \label{sm:offset-tip}
\end{equation}
in endpoint translation and rotation coordinates. The benchmark sets
$e=0$. Coating, adhesive and mounting masses are not represented by
this ideal bare-core inertia.

The nodal displacement entries have units of length and the slope
entries are dimensionless. Consequently, the $(w,w)$, $(w,\vartheta)$
and $(\vartheta,\vartheta)$ blocks of $M_f$ have units kg, kg\,m and
kg\,m$^2$, respectively. The corresponding stiffness units are
N/m, N and N\,m; the generalized force entries are N and N\,m.
These mixed units explain why neither a scalar mass replacement nor
an unscaled Euclidean actuator condition number is physically neutral.

\subsection{Linear frequency verification and parameter sensitivity}
\label{sm:fe-check}
An independent continuous-beam check is available because the biased
linearized endpoint problem has constant coefficients. Put
\begin{align}
 \lambda&=kL,\qquad
 \omega^2=\frac{EI}{\rho A_cL^4}\lambda^4,\nonumber\\
 w(s)&=A[\cosh(ks)-\cos(ks)]\nonumber\\
 &\quad+B[\sinh(ks)-\sin(ks)].
 \label{sm:characteristic-shape}
\end{align}
The undamped endpoint conditions for this eigenproblem are
\begin{align}
 EIw_{sss}(L)+m_t\omega^2w(L)&=0,\nonumber\\
 EIw_{ss}(L)+(\mu B_0-J_t\omega^2)w_s(L)&=0.
 \label{sm:characteristic-bc}
\end{align}
Define $\eta_t=m_t/(\rho A_cL)$,
$\iota_t=J_t/(\rho A_cL^3)$,
$\kappa_t=\mu B_0L/(EI)$ and
\begin{align*}
 f&=\cosh\lambda-\cos\lambda,&
 g_\lambda&=\sinh\lambda-\sin\lambda,\\
 c_+&=\cosh\lambda+\cos\lambda,&
 s_+&=\sinh\lambda+\sin\lambda.
\end{align*}
Writing $h_\lambda=\kappa_t/\lambda-\iota_t\lambda^3$, the positive
roots satisfy
\begin{equation}
 \det\begin{pmatrix}
 g_\lambda+\eta_t\lambda f&
 c_++\eta_t\lambda g_\lambda\\
 c_++h_\lambda s_+&
 s_++h_\lambda f
 \end{pmatrix}=0.
 \label{sm:characteristic-determinant}
\end{equation}
For zero tip inertia and zero bias this reduces to
$1+\cosh\lambda\cos\lambda=0$. The damped frequencies are not used in
mass normalization; $\omega_j$ throughout the paper denotes the
undamped frequency of the biased stiffness.

\begin{table*}[t]
\centering\small
\caption{Angular frequencies in rad/s for the revised beam, including
magnetic-core translational and rotational inertia. The final row solves
the independent continuous-beam endpoint problem.}
\label{sm:frequency-table}
\begin{tabular}{lrrrr}
\toprule
Discretization & $\omega_1$ & $\omega_2$ & $\omega_3$ & $\omega_4$\\
\midrule
$16$ Hermite elements & $7.840010222$ & $47.399809441$ & $131.211581398$ & $256.375395140$\\
$32$ Hermite elements & $7.840009162$ & $47.399576624$ & $131.206681073$ & $256.339221935$\\
$64$ Hermite elements & $7.840009096$ & $47.399561992$ & $131.206371332$ & $256.336917925$\\
Continuous endpoint equation & $7.840009091$ & $47.399561016$ & $131.206350623$ & $256.336763590$\\
\bottomrule
\end{tabular}
\end{table*}

The eigenvectors are mass-normalized and oriented so that each tip
displacement coefficient is positive. The $32$-element maximum
mass-orthonormality defect is $5.14\times10^{-15}$; the first two
relative frequency errors against \eqref{sm:characteristic-determinant}
are $9.02\times10^{-9}$ and $3.29\times10^{-7}$.
At fixed bias, omitting tip inertia gives
$(\omega_1,\omega_2)=(7.859186509,47.517207103)$ rad/s.
Including it lowers the frequencies by approximately $0.2440\%$ and
$0.2476\%$. The linear verification does not bound errors caused by
neglecting nonlinear high-mode response.

Table~\ref{sm:sensitivity-table} varies parameters one at a time.
The $B_r$ endpoints come from the N50 grade interval \citep{smN50}.
The inertia and modulus perturbations are declared design studies,
not measured confidence intervals. Inertia perturbations scale $m_t$
and $J_t$ together at fixed core dimension and fixed magnetic moment.
For reported damping-ratio changes in the machine-readable record,
the nominal $\alpha$ is held fixed.

\begin{table*}[t]
\centering\small
\caption{Design sensitivity of the first two biased frequencies.
Percentages are relative to the nominal $32$-element model.}
\label{sm:sensitivity-table}
\begin{tabular}{lrrl}
\toprule
Change & $\Delta\omega_1/\omega_1$ & $\Delta\omega_2/\omega_2$ & Interpretation\\
\midrule
$B_r=1.40$ T & $-0.082935\%$ & $-0.029460\%$ & Grade-range endpoint\\
$B_r=1.45$ T & $+0.082656\%$ & $+0.029413\%$ & Grade-range endpoint\\
$m_t,J_t$ reduced by $10\%$ & $+0.024384\%$ & $+0.024646\%$ & Chosen design perturbation\\
$m_t,J_t$ increased by $10\%$ & $-0.024367\%$ & $-0.024608\%$ & Chosen design perturbation\\
$E$ reduced by $20\%$ & $-9.526743\%$ & $-10.186350\%$ & Chosen design perturbation\\
$E$ increased by $20\%$ & $+8.668840\%$ & $+9.235817\%$ & Chosen design perturbation\\
\bottomrule
\end{tabular}
\end{table*}

The combined $B_r$ and $\pm10\%$ inertia corners change $\omega_1$ by
between $-0.10729\%$ and $+0.10705\%$. The tested modulus changes
therefore dominate those magnetic-core perturbations in this frequency
comparison. They do not establish a bound on all possible fabrication
variation.

The material and density reference is a particular Sylgard preparation
\citep{cho2021pdms}. Mechanical loss depends on preparation and
measurement conditions; loss-factor measurements for other samples
\citep{fischer2017adhesion} are not converted into the present modal
damping ratios without a constitutive and frequency correspondence.
The scenarios $\zeta_1^{\mathrm{d}}\in\{0,0.01,0.025,0.05,0.1\}$ instead define
explicit effective damping models. Likewise, the core dimensions are
design inputs and the remanent-field midpoint is nominal; the
available specifications do not calibrate a coated, bonded specimen.

\subsection{Canonical units, complex derivatives and input conversion}
\label{sm:scaling}
After mass normalization, $a_j$ has units
$\sqrt{\mathrm{kg}}\,\mathrm m$ and
$\pi_j=\dot a_j$ has units $\sqrt{\mathrm{kg}}\,\mathrm{m\,s^{-1}}$.
Let $P_s=\operatorname{diag}(\sqrt{E_0\omega_j/\omega_0})$.
The scales in \eqref{mech:scaling} satisfy
\begin{align}
 D_sP_s&=(E_0/\omega_0)I,\nonumber\\
 \sum_j\dd a_j\wedge\dd\pi_j
 &=\frac{E_0}{\omega_0}\sum_j\dd X_j\wedge\dd Y_j.
 \label{sm:scaled-form}
\end{align}
With $h=H/E_0$ and $\tau=\omega_0t$, the chain rule gives
$X'=h_Y$, $Y'=-h_X$. The physical quadratic kinetic energy becomes
$E_0Y^{\mathsf T}\mathsf NY/2$, while the Lagrangian kinetic energy
in dimensionless position and velocity is
$v^{\mathsf T}\mathsf N^{-1}v/2$.
The three distinct masses are therefore the nodal $M_f$, the
mass-normalized modal identity and the balanced
$M_{\rm eff}=\mathsf N^{-1}$.

The physical Rayleigh potential
$\mathcal D_{\rm phys}=\alpha\dot a^{\mathsf T}\dot a/2$
transforms as
\begin{equation}
 \mathcal D(X,v)=
 \frac{\mathcal D_{\rm phys}(D_sX,\omega_0D_sv)}{E_0\omega_0}
 =\tfrac12\gamma v^{\mathsf T}\mathsf N^{-1}v.
 \label{sm:rayleigh-scale}
\end{equation}
For the fixed kinetic energy, $v=\mathsf NY$ and its force is
$-\partial_v\mathcal D=-\gamma Y$. If learned momentum terms change
$h_Y$, the force remains $-\gamma\mathsf N^{-1}h_Y$, rather than
automatically $-\gamma Y$. The realization in
\eqref{app:full-feedback} handles a general target through its actual
velocity map.

For $\zeta=(X+\mathrm iY)/\sqrt2$ and
$\chi=(X-\mathrm iY)/\sqrt2$, the inverse chart gives
\begin{equation}
 \partial_\chi h^{\mathbb C}
     =(\partial_Xh+\mathrm i\partial_Yh)/\sqrt2.
 \label{sm:complex-chain-rule}
\end{equation}
Substituting the canonical equations into $\zeta'$ proves
\eqref{mech:wirtinger} with no additional factor of two.
For a force $Q$ acting in the $Y$ equation, the extra term in
$\zeta'$ is $\mathrm iQ/\sqrt2$. A proposed complex readout
$h+\mathrm i\mathcal D$ must therefore be used with the decoder
\eqref{mech:decoder}; differentiating that sum as though it were one
conservative Hamiltonian generally produces another vector field.

For two modes, $b_j$ has units $\mathrm{kg}^{-1/2}$ and
$t_j$ has units $\mathrm{kg}^{-1/2}\mathrm m^{-1}$.
Both $b_jF$ and $t_j\tau_B$ consequently have the units of modal
acceleration. From $X'=\mathsf NY$,
\begin{equation}
 \ddot a=\omega_0^2D_s\mathsf NY'=\Lambda Y',\qquad
 \Lambda=E_0D_s^{-1}.
 \label{sm:force-scale}
\end{equation}
Moreover, $d=D_s\boldsymbol t$ gives
\begin{equation}
 \Lambda\nabla V_0
 =\Omega_2^2a+\boldsymbol t E_0(\sin\vartheta-\vartheta).
 \label{sm:scaled-potential-force}
\end{equation}
Equations \eqref{sm:force-scale}--\eqref{sm:scaled-potential-force}
prove \eqref{app:residual-plant}. If a program stores the total modal
moment $T$, it must first remove
$E_0(\vartheta-\sin\vartheta)$ before dividing by
$\mu\cos\vartheta$ to obtain $B_y$.
This subtraction changes a port convention; it does not compensate the
unknown residual $-\Lambda\nabla R_*$.
For the learned input, $\Lambda\widehat\gamma Y=
\omega_0\widehat\gamma\dot a=\widehat\alpha\dot a$.
The physical potential forces combine as
$-\Lambda\nabla R_*+\Lambda\nabla\widehat R
=-\Lambda\nabla\delta R$, which completes
\eqref{app:physical-closed-loop}. No true residual derivative is used
by the fitted controller; it is evaluated only by the synthetic plant
or its explicitly labeled truth-only comparator.


The numerical conversion uses
\begin{align}
 A_{\rm mag}&\simeq
 \begin{pmatrix}
 203.96232682&3303.68875398\\
 205.00706618&11909.57162007
 \end{pmatrix},\nonumber\\
 D_s&=\operatorname{diag}(S_1,S_2),\nonumber\\
 S_1&\simeq9.29943515\times10^{-5},\nonumber\\
 S_2&\simeq3.78205312\times10^{-5}.
 \label{sm:numerical-ports}
\end{align}
Its determinant is nonzero. Rescaling the columns by actual available
force and torque is needed to assess a device's effective authority;
column normalization alone verifies only geometric separation of the
two generalized directions.

\subsection{Residual compensation under modal enrichment}
\label{sm:enriched}
The full retained-mode potential corresponding to
\eqref{app:enriched-plant} is
\begin{align}
 V_N&=\tfrac12a^{(N)\mathsf T}\Omega_N^2a^{(N)}
       +E_0g(\vartheta_N)\nonumber\\
 &\quad+\mathcal R_N(a^{(N)}).
 \label{sm:enriched-potential}
\end{align}
Its kinetic energy is $\dot a^{(N)\mathsf T}\dot a^{(N)}/2$.
The quadratic magnetic bias is already in $\Omega_N^2$, so the same
subtraction in $g$ avoids double counting. All potentials
$\mathcal R_N$ in a comparison are restrictions of one specified
full potential. The compatibility condition in the main text ensures
that setting the high coordinates to zero recovers the original
synthetic residual, without asserting that its continuation was
identified experimentally.

For the first two modal equations,
\begin{align}
 \Lambda^{-1}\Omega_2^2a_{1:2}&=\mathsf N X,\nonumber\\
 \Lambda^{-1}E_0\boldsymbol t_{1:2}&=d,\nonumber\\
 \Lambda^{-1}\nabla_{a_{1:2}}\mathcal R_N
 &=\nabla R_*(X)+e_R^{(N)}.
 \label{sm:enriched-projection}
\end{align}
Dividing their acceleration equation by $\Lambda$ and inserting
$U_{\rm rec}=\widehat\gamma Y+\nabla\widehat R+e_{\rm port}$
proves \eqref{app:enriched-closed-loop}. In particular,
$\vartheta_N-\vartheta_2=\sum_{j=3}^Nt_ja_j$.
If both angles lie in $[-\vartheta_{\max},\vartheta_{\max}]$, the
mean-value theorem and $1-\cos s\le s^2/2$ give
\begin{equation}
 \|e_{\rm mag}^{(N)}\|
 \le\tfrac12\|d\|\vartheta_{\max}^2
                   |\vartheta_N-\vartheta_2|.
 \label{sm:angle-defect-bound}
\end{equation}
A bound on high-mode motion is still needed to use this inequality.
Sampling small angles on a two-mode trajectory does not supply it.

For the declared residual $R_*(X)=\beta_*g(\vartheta_2)$, one concrete
synthetic extension for the planned modal comparison is
\begin{equation}
 \mathcal R_N^{\rm ang}(a^{(N)})=E_0\beta_*g(\vartheta_N).
 \label{sm:angle-residual-extension}
\end{equation}
It preserves the low-dimensional restriction and gives
$e_R^{(N)}=-\beta_*e_{\rm mag}^{(N)}$. Thus the combined magnetic
and residual difference is $(1+\beta_*)e_{\rm mag}^{(N)}$, even with
perfect coefficient identification. As an isolating comparison, the
alternative
$\mathcal R_N^{\rm low}=E_0R_*(D_s^{-1}a_{1:2})$ gives
$e_R^{(N)}=0$ and no residual force in the higher coordinates.
These are different declared synthetic potentials, not equivalent
consequences of the same two-mode observations. Neither higher-mode
comparison has been executed in the present numerical record.

Let $H_{0,N}$ omit $\mathcal R_N$ from the full retained energy.
For the applied physical ports,
\begin{align}
 \dot H_{0,N}
 &=-\alpha\|\dot a^{(N)}\|^2
       -\dot a^{(N)\mathsf T}\nabla_a\mathcal R_N\nonumber\\
 &\quad+F_{\rm app}\dot w(L)+\tau_{B,\rm app}\dot\vartheta_N.
 \label{sm:enriched-power}
\end{align}
The power contributed by the residual must be included in a
nominal-energy work bound. Alternatively, the augmented energy may be used after a
suitable lower bound by bending energy has been established. Its
nonnegative decomposition cannot be assumed from the two-mode
residual ansatz.

\subsection{A stronger exact realization of a conservative target}
\label{sm:exact-lift}
This construction uses known damping and exact target derivatives.
It is retained as a conditional full-model result separate from the
unchanged fitted feedback in \eqref{app:enriched-closed-loop}.
For the nominal finite-element plant \eqref{mech:fe-plant}, retain all
$N_f=64$ modes and define
\begin{align}
 \ddot a_j+\alpha\dot a_j+\omega_j^2a_j&=b_jF+t_jT,
       \qquad j=1,\ldots,N_f,\nonumber\\
 T&=E_0(\vartheta-\sin\vartheta)+\tau_B,\nonumber\\
 \vartheta&=\sum_{j=1}^{N_f}t_ja_j.
 \label{app:full-modal}
\end{align}
The symbol $T$ is a total modal moment, whereas $U$ remains the
dimensionless generalized input. Equation~\eqref{app:full-modal}
specifies the full plant for the following result; an arbitrary
unmodelled force on omitted modes is not included in its hypotheses.

Consider a real analytic target Hamiltonian $h_{\rm tar}(X,Y)$, which may
contain momentum-dependent terms. With physical low-mode velocity
$v_{1:2}=\dot a_{1:2}$, infer $Y$ from
\begin{equation}
 (h_{\rm tar})_Y(X,Y)=\omega_0^{-1}D_s^{-1}v_{1:2},
 \qquad X=D_s^{-1}a_{1:2}.
 \label{app:velocity-inversion}
\end{equation}
On a domain with positive definite $(h_{\rm tar})_{YY}$ and convex
$Y$-fibres this relation is injective in $Y$ and locally invertible.
The acceleration and port requests that reproduce the canonical target
are
\begin{align}
 a^{\rm des}_{1:2,tt}
 &=D_s\omega_0^2\bigl[
     (h_{\rm tar})_{YX}(h_{\rm tar})_Y
       -(h_{\rm tar})_{YY}(h_{\rm tar})_X\bigr],\nonumber\\
 \binom{F}{T}
 &=A_{\rm mag}^{-1}
   \bigl[a^{\rm des}_{1:2,tt}+\alpha v_{1:2}
                        +\Omega_2^2a_{1:2}\bigr],\nonumber\\
 B_y&=\frac{T-E_0(\vartheta-\sin\vartheta)}{\mu\cos\vartheta}.
 \label{app:full-feedback}
\end{align}
The actual full tip angle enters the last line. Substitution first
enforces the desired low acceleration; differentiating
\eqref{app:velocity-inversion} then recovers both canonical equations.

\begin{proposition}[Conditional finite-element lift]
\label{app:conditional-lift}
Suppose $h_{\rm tar}$ has an analytic invariant embedded two-torus
carrying a quasiperiodic flow, and the
velocity map in \eqref{app:velocity-inversion} satisfies the stated
positivity and fibre conditions on a neighbourhood of it. Assume
$\alpha>0$, $A_{\rm mag}$ is invertible and the exact feedback
\eqref{app:full-feedback} is available. If the resulting response graph
lies inside the allowed angle and actuator domain, it is an invariant
two-torus of \eqref{app:full-modal}. For each fixed low-mode orbit on
this torus, higher-mode deviations decay exponentially while the
closed-loop trajectory remains in that domain.
\end{proposition}


The feedback in \eqref{app:full-feedback} cancels the nominal full-angle
remainder to impose a chosen low acceleration. It does not contain an
arbitrary unknown full residual $\mathcal R_N$, and it is not obtained
merely by applying the learned two-mode law to more modes. The result
concerns a finite-dimensional plant; no infinite-dimensional beam
limit or attraction of arbitrary low-mode initial conditions is
asserted. The proof and its domain requirements follow below.

\subsection{Proof and scope of the finite-element lift}
\label{sm:lift}
The proof uses the nominal full modal plant
\eqref{app:full-modal}, known parameters and exact feedback.
Let $\iota(\phi)=(X(\phi),Y(\phi))$ be the target torus, with
$\phi\in(\mathbb R/2\pi\mathbb Z)^2$ and physical frequency $\Omega$.
Its physical low-mode image is
\begin{equation}
 a_{1:2}=D_sX(\phi),\qquad
 v_{1:2}=\omega_0D_s(h_{\rm tar})_Y(\iota(\phi)).
 \label{sm:physical-embedding}
\end{equation}
If two points in this image have the same position, they have the same
$X$. For distinct $Y_1,Y_2$ in the same convex fibre, write
$\Delta Y=Y_1-Y_2$, $Y_s=Y_2+s\Delta Y$ and $h=h_{\rm tar}$. Then
\begin{align}
 &\Delta Y^{\mathsf T}[h_Y(X,Y_1)-h_Y(X,Y_2)]\nonumber\\
 &\quad=\int_0^1
 \Delta Y^{\mathsf T}h_{YY}(X,Y_s)\Delta Y\dd s>0.
 \label{sm:velocity-injectivity}
\end{align}
Thus equality of the physical velocity also implies $Y_1=Y_2$.
The differential of the velocity map is invertible in the $Y$
directions, so it preserves the embedded low torus locally and
globally on the stated domain.

The requested $F,T$ in \eqref{app:full-feedback} depend only on the low
physical state. The full tip angle enters solely in conversion from
$T$ to $B_y$, which cancels the same nonlinear bias remainder in every
modal equation. Each higher mode then has prescribed quasiperiodic
forcing. For a real temporal frequency $\sigma$,
\begin{equation}
 D_j(\sigma)=\omega_j^2-\sigma^2+\mathrm i\alpha\sigma
 \label{sm:stable-resolvent}
\end{equation}
cannot vanish: at $\sigma=0$ its real part is positive, and at nonzero
$\sigma$ its imaginary part is nonzero. Its modulus tends to infinity
as $|\sigma|\to\infty$, so its reciprocal is bounded on the real
axis. Using the Fourier convention $e^{\mathrm i k\cdot\phi}$,
the unique bounded higher-mode response therefore has coefficients
\begin{equation}
 a^*_{j,k}=
 \frac{b_jF_k+t_jT_k}
 {\omega_j^2-(k\cdot\Omega)^2+\mathrm i\alpha(k\cdot\Omega)},
 \qquad j\ge3.
 \label{app:response-coefficients}
\end{equation}
Consequently the exponentially decaying Fourier coefficients of
$F,T$ remain exponentially decaying after multiplication by
$D_j(k\cdot\Omega)^{-1}$. This establishes an analytic response
for each of the finitely many higher modes.

Differences from that response satisfy
$\ddot z_j+\alpha\dot z_j+\omega_j^2z_j=0$ and decay exponentially.
The low motion remains neutral along its target torus; no attraction
of general low initial conditions follows. With $\alpha=0$, the
resolvent may have resonances and this particular lifting argument
does not apply. That distinction includes the zero-damping sensitivity
scenario.

The result requires the complete response graph and any considered
startup transient to remain in the angle, sensing and actuator domain.
It does not supply those bounds merely from the two-mode fit.
In particular, the reduced synthetic potential residual is not a
specified full-space model of omitted-mode forces. Extending the
theorem to such forces requires an explicit full model and its
corresponding estimates. Similarly, a governor that changes $F,T$
breaks exact target realization even when it preserves clearance.

\subsection{Whole-centreline energy bound and initialization}
\label{sm:clearance}
For $w\in H^2(0,L)$ with $w(0)=w_s(0)=0$,
\begin{align}
 |w(s)|^2
 &\le\left(\int_0^s(s-\xi)^2\dd\xi\right)
                \left(\int_0^s|w_{ss}|^2\dd\xi\right)\nonumber\\
 &\le\frac{2s^3}{3EI}E_{\rm bend},\qquad
 E_{\rm bend}=\frac{EI}{2}\int_0^L|w_{ss}|^2\dd s.
 \label{sm:clearance-proof}
\end{align}
Since the kinetic energy includes a positive tip mass and inertia,
and $E_0(1-\cos\vartheta)\ge0$, the nominal total energy satisfies
$E\ge E_{\rm bend}$. Taking the supremum proves
\eqref{app:energy-bound}. This also holds for conforming $C^1$
Hermite displacement fields. The related slope estimate is
\begin{equation}
 \|w_s\|_\infty^2
 \le L\int_0^L|w_{ss}|^2\dd s
 \le\frac{2LE}{EI}.
 \label{sm:slope-energy}
\end{equation}
These inequalities require a bound on the relevant total physical
energy, not just on one modal amplitude or the tip displacement.

The constant in \eqref{sm:clearance-proof} is sharp for the bending
energy at the tip: equality is attained by a curvature proportional
to $L-s$, the static tip-force shape. That equality case explains
the familiar expression $E_c=3EIc^2/(2L^3)$ but does not restrict the
inequality to static shapes.
For the revised beam $E_c=10.0197171$ nJ.
An energy cap $E_{\max}=0.95E_c=9.5187312$ nJ would imply
$\|w\|_\infty\le\sqrt{0.95}\,c<c$.

An initialization bound can use the full-state quadratic form
\begin{equation}
 Q(q,v)=\tfrac12v^{\mathsf T}M_fv+
        \tfrac12q^{\mathsf T}K_bq.
 \label{sm:initial-quadratic}
\end{equation}
Since $1-\cos\vartheta\le\vartheta^2/2$, $E(q,v)\le Q(q,v)$.
If estimates $(\widehat q,\widehat v)$ satisfy
$Q(q-\widehat q,v-\widehat v)\le\mathcal E_{\rm init}$, then
\begin{equation}
 E(q,v)\le
 \left[\sqrt{Q(\widehat q,\widehat v)}
                  +\sqrt{\mathcal E_{\rm init}}\right]^2.
 \label{sm:initial-upper-bound}
\end{equation}
This follows from the triangle inequality for $\sqrt Q$.
Tip-only measurements do not by themselves bound energy in unobserved
bending or velocity components.

For the stored nominal two-mode initial state
$(X,Y)=(0.09,0.08,0.015,-0.01)$,
direct substitution into $E_0h_0$ gives $12.6567207$ nJ.
This exceeds $E_{\max}$, so the work-governor initialization condition
has not been verified for the reported compensation runs. Their
sampled centreline values can be below the comparison clearance
without satisfying this sufficient energy condition. The tracking
initialization belongs to the separate extension in
Supplementary Section~\ref{sec:supp-tracking}.

For a plant augmented by $R_*$, the nominal-energy derivative gains
the power of its additional generalized force. Alternatively, one
could use the augmented energy only after proving a suitable lower
bound by bending energy. The synthetic residual alone is therefore
not silently incorporated into the nominal clearance theorem.

\subsection{Input-work governor and bounded implementation errors}
\label{sm:work}
For the nominal power identity and exact port measurements,
\begin{equation}
 E(t)-W(t)=E(0)-W(0)
             -\int_0^t\alpha\dot q^{\mathsf T}M_f\dot q\dd s\le0.
 \label{sm:work-ledger}
\end{equation}
The constraint in \eqref{app:work-governor} implies
$\dot W\le\kappa(E_{\max}-W)$.
With $W(0)<E_{\max}$, comparison gives
$E_{\max}-W(t)\ge[E_{\max}-W(0)]e^{-\kappa t}$.
The governor has a feasible zero active input whenever the prescribed
bias itself is feasible. For a nonempty closed convex admissible port set containing zero and a
positive definite scaled cost, its quadratic projection is unique.

The same construction can accommodate deterministic error envelopes,
provided they are supplied independently. Let the applied ports obey
\begin{align}
 F&=(1+g_F)u_F+e_F,\nonumber\\
 B_y&=(1+g_B)u_B+e_B,\nonumber\\
 |g_i|&\le\eta_i,\qquad |e_i|\le\epsilon_i.
 \label{sm:port-errors}
\end{align}
For errors bounded by
$|v_w-\widehat v_w|\le\delta_v$,
$|\vartheta-\widehat\vartheta|\le\delta_\vartheta$ and
$|\dot\vartheta-\widehat{\dot\vartheta}|\le\delta_\omega$, define
\begin{align}
 \widehat a_F&=\widehat v_w,\qquad &&\Delta_F=\delta_v,\nonumber\\
 \widehat a_B&=\mu\cos\widehat\vartheta\,\widehat{\dot\vartheta},
 &&\Delta_B=\mu(\delta_\omega+
                   |\widehat{\dot\vartheta}|\delta_\vartheta).
 \label{sm:power-estimation}
\end{align}
The cosine is one-Lipschitz, so $|a_i-\widehat a_i|\le\Delta_i$.
A convex positively homogeneous upper bound on the non-additive port
power is
\begin{equation}
 P^+(u)=\sum_{i\in\{F,B\}}
 \left[\widehat a_i u_i+
   \{\Delta_i+\eta_i(|\widehat a_i|+\Delta_i)\}|u_i|\right].
 \label{sm:power-envelope}
\end{equation}
It contains no division by a measured velocity and satisfies
$P^+(0)=0$.

For $\alpha>0$, define inverse-mass coefficients
$m_F=e_w^{\mathsf T}M_f^{-1}e_w$ and
$m_\vartheta=e_\vartheta^{\mathsf T}M_f^{-1}e_\vartheta$.
The additive generalized force
$r_e=e_we_F+\mu e_\vartheta e_B\cos\vartheta$ obeys
\begin{align}
 r_e^{\mathsf T}v-\alpha v^{\mathsf T}M_fv
 &\le\frac{r_e^{\mathsf T}M_f^{-1}r_e}{4\alpha}\le p_e,\nonumber\\
 p_e&=\frac{(\epsilon_F\sqrt{m_F}
                  +\mu\epsilon_B\sqrt{m_\vartheta})^2}{4\alpha}.
 \label{sm:additive-power}
\end{align}
Completion of the square proves the first inequality, and the triangle
inequality in the $M_f^{-1}$ norm proves the second.
Thus $\dot E\le P^+(u)+p_e$.

On a chosen horizon $[0,T_H]$, reserve $R_H=T_Hp_e$ and put
$C_H=E_{\max}-R_H$. If $W(0)\ge E(0)$ and
$W(0)<C_H$, integrate $\dot W=P^+(u)$ and impose
\begin{equation}
 P^+(u)\le\kappa(C_H-W),\qquad u\in\mathcal U(q).
 \label{sm:robust-governor}
\end{equation}
Then, throughout $0\le t\le T_H$,
\begin{equation}
 E(t)\le W(t)+tp_e<C_H+R_H=E_{\max}.
 \label{sm:robust-energy}
\end{equation}
The command set must reserve the stated gain and additive errors when
enforcing actual actuator limits. If additive errors vanish, the
reserve is zero and the exact-power argument extends to all time for
$\alpha\ge0$. A nonzero constant additive-error allowance gives the
stated finite-horizon condition, not an all-time conclusion.

For sample-and-hold execution, a bound at one sampling instant is
insufficient. If $P_k^+(u_k)$ bounds power throughout the hold interval
of length $h_k$, the rule
\begin{equation}
 h_kP_k^+(u_k)\le C_H-W_k,\qquad
 W_{k+1}=W_k+h_kP_k^+(u_k)
 \label{sm:discrete-work}
\end{equation}
also bounds the energy within that interval. The error envelope must
cover sensor age, hold age and command latency. For example, a known
command slew bound $|\dot u_i|\le S_i$ and delay $d_i$ may be included
by increasing $\epsilon_i$ by $(1+\eta_i)S_id_i$.
These are conditional implementation bounds; no numerical sensor,
latency or coil-current envelope has been calibrated for the revised
benchmark. The reported runs test the ideal physical input conversion,
while this work argument states what an energy-certified implementation
would additionally have to satisfy.
